
\textbf{\Huge\newline\newline\newline\newline\newline Resumen}


\vspace{2cm}
\bigskip

Pendiente...

%Este proyecto consiste en la creación de un prototipo de una estación de predicción de fallas en unidades de ''balloon'' de catéteres cardiacos, por medio de procesamiento de imágenes utilizando una red neuronal convolucional, la cual aprende las características de las unidades por medio de fotografías tomadas de las unidades mientras se les aplica luz polarizada. El sistema completo consiste de una interfaz de usuario que le permite al operador ingresar sus credenciales, identificar las unidades, realizar la captura e importación de fotografías y realizar el análisis de forma automatizada, así cómo un hardware que permite sostener y rotas las unidades frente a una lampara de luz polarizada y tomar la fotografías de la unidad posicionada al presionar un botón en la interfaz de usuario del software.

%Para el desarrollo de este proyecto se utilizó el lenguaje de programación Python 3 y la biblioteca Tensorflow (más específicamente el API de Keras incluido en este) para la creación del modelo de la red neuronal, su entrenamiento y su implementación. Se utilizó la biblioteca pyQT5 para la creación de la interfaz de usuario, la biblioteca OpenCV para el control de la cámara que captura las imágenes y el almacenamiento en disco duro de estas. Se hizo uso de una plataforma Arduino UNO R3 para realizar el control de un motor paso a paso 28BYJ-48 por medio de un módulo controlador ULN2003, así cómo el envío y recibimiento de comandos por medio de su puerto UART.
%Para la captura de la fotografías se utiliza un microscopio de mano Celestron Capture Pro 5MP compatible con USB.


%En este documento se detalla el proceso de diseño, creación, prueba y validación del sistema, los resultados obtenidos y las diferentes partes que lo componen. Cabe resaltar que detalles considerados cómo confidenciales por parte de la empresa serán omitidos por solicitud de la misma.


\textbf{\newline Palabras clave: A, B, C, D, E} 

%\textbf{\newline Palabras clave: Arduino, Celestron, Redes neuronales convolucionales, Procesamiento de imágenes, Luz polarizada, Visión por computadora, Catéter cardiaco, Tensorflow, Keras, OpenCV} 



\newpage
\textbf{\Huge\newline\newline\newline\newline\newline Abstract}
\vspace{2cm}
\bigskip

Pending...

%This project consists in the creation of a prototype station for failure prediction on balloon units for cardiac catheters using image processing, employing a convolutional neural network which learns the features of the units by pictures taken of these units while applying through them polarized light. The complete system consists of an user interface that allows the operator to input his credentials, identify the units, capture and import pictures and analyze the pictures automatically. Also, a hardware that allows to hold and rotate the units in front of the polarized light lamp and take the pictures of the positioned unit when a button in the user interface of the software is pressed.

%For the development of this project the Python programming language and the Tensorflow library (the Keras API that it includes more specifically) were used for the creation of the neural network's model, it's training and implementation. PyQT5 library was used for the development of the user interface, OpenCV library for the camera control to take the pictures and disc storage of these. An Arduino UNO R3 platform was employed to control a 28BYJ-48 stepper motor using a ULN2003 controller, also the send and receiving of commands using the Arduino's UART port.
%For the picture capture an USB compatible handheld Celestron Capture Pro 5MP microscope was used.
 
%This document details the design, creation, testing and validation process of the system, also the obtained results and the different parts that make it up. It fits to highlight that details considered as confidential by the company will be omitted as per it's request.


\textbf{\newline Keywords: A, B, C, D, E} 


%\textbf{\newline Keywords: Arduino, Celestron, Convolutional Neural Networks, Image Processing, Polarized Light, Computer Vision, Cardiac Catheter, Tensorflow, Keras, OpenCV} 


%%% Local Variables: 
%%% mode: latex
%%% TeX-master: "main"
%%% End: 
